\section{随机变量的数学期望（均值）}
\textbf{离散型 r.v. 的期望}：
\begin{equation}
	E(X) = \sum_{k=1}^{\infty} x_k p_k \quad \left(\text{要求 } \sum_{k=1}^{\infty} |x_k|\,p_k < \infty\right)
	\label{eq:第四章:离散期望}
\end{equation}
\textbf{连续型 r.v. 的期望}：
\begin{equation}
	E(X) = \int_{-\infty}^{\infty} x\,f(x)\,dx \quad \left(\text{要求 } \int_{-\infty}^{\infty} |x|\,f(x)\,dx < \infty\right)
	\label{eq:第四章:连续期望}
\end{equation}
\textbf{函数期望}：
\begin{itemize}
	\item \textbf{一维}：$E[g(X)] = \int_{-\infty}^{\infty} g(x) f(x)\,dx$（连续型）
	\item \textbf{二维}：$E[g(X, Y)] = \int_{-\infty}^{\infty} \int_{-\infty}^{\infty} g(x, y) f(x, y)\,dx\,dy$
\end{itemize}

\subsection{期望的性质}
\textbf{1. 常数的期望}：如果 $C$ 是一个常数，那么它的期望就是它本身。
\begin{equation}
	E(C) = C
	\label{eq:第四章:常数期望}
\end{equation}
理解：常数没有波动，其平均值自然就是它自己。

\textbf{2. 线性性质（Linearity of Expectation）}：这是期望最核心、最常用的性质，分为两部分：
\begin{itemize}
	\item 倍数关系：随机变量乘以常数 $a$，期望也扩大 $a$ 倍：
	\begin{equation}
		E(aX) = aE(X)
		\label{eq:第四章:倍数}
	\end{equation}
	\item 和的期望：两个随机变量之和的期望等于各自期望之和：
	\begin{equation}
		E(X + Y) = E(X) + E(Y)
		\label{eq:第四章:和的期望}
	\end{equation}
	注意：这一条对于任意两个随机变量都成立，无论它们是否独立。
\end{itemize}
将上述结合，可以得到更通用的线性公式：
\begin{equation}
	E(aX + bY + c) = aE(X) + bE(Y) + c
	\label{eq:第四章:线性组合}
\end{equation}

\textbf{3. 独立随机变量的乘积}：如果随机变量 $X$ 和 $Y$ 是相互独立的，那么它们乘积的期望等于各自期望的乘积：
\begin{equation}
	E(XY) = E(X)E(Y)
	\label{eq:第四章:乘积期望}
\end{equation}
警示：如果 $X$ 和 $Y$ 不独立，这个等式通常不成立。

\textbf{4. 函数的期望（EUE 定理）}：若 $Y$ 是 $X$ 的函数，即 $Y = g(X)$，我们不需要求出 $Y$ 的分布，可以直接利用 $X$ 的概率密度或分布律求期望：
\begin{itemize}
	\item \textbf{离散型}：$E[g(X)] = \sum_{i} g(x_i)\,P(X = x_i)$
	\item \textbf{连续型}：$E[g(X)] = \int_{-\infty}^{+\infty} g(x) f(x)\,dx$
\end{itemize}

\textbf{5. 单调性}：如果对于所有的取值，都有 $X \le Y$，那么它们的期望也满足：
\begin{equation}
	E(X) \le E(Y)
	\label{eq:第四章:单调性}
\end{equation}
由此推论：如果 $X \ge 0$，则 $E(X) \ge 0$。

\subsubsection{常见随机变量的期望参考表}
\begin{table}[H]
\centering
\begin{tabular}{lcc}
\toprule
常用分布 & 符号 & 期望 $E(X)$ \\
\midrule
0-1 分布 & $B(1, p)$ & $p$ \\
二项分布 & $B(n, p)$ & $np$ \\
泊松分布 & $P(\lambda)$ & $\lambda$ \\
均匀分布 & $U(a, b)$ & $\dfrac{a+b}{2}$ \\
正态分布 & $N(\mu, \sigma^2)$ & $\mu$ \\
指数分布 & $e(\lambda)$ & $\dfrac{1}{\lambda}$ \\
\bottomrule
\end{tabular}
\end{table}

\section{方差}
\textbf{定义}：
\begin{equation}
	D(X) = E\big[(X - E(X))^{2}\big]
	\label{eq:第四章:方差定义}
\end{equation}
\textbf{计算公式}：
\begin{equation}
	D(X) = E(X^{2}) - [E(X)]^{2}
	\label{eq:第四章:方差计算}
\end{equation}
\textbf{性质}：
\begin{enumerate}
	\item $D(c) = 0$；若 $D(X) = 0$，则 $P\{X = c\} = 1$。
	\item $D(aX + b) = a^{2} D(X)$。
	\item 若 $X, Y$ 独立，则 $D(X + Y) = D(X) + D(Y)$。
\end{enumerate}
\textbf{重要 r.v. 的方差}：
\begin{enumerate}
	\item \textbf{二项分布 $b(n, p)$}：$D(X) = n p (1-p)$
	\item \textbf{指数分布 $E(\lambda)$}：$D(X) = \frac{1}{\lambda^{2}}$
	\item \textbf{均匀分布 $U(a, b)$}：$D(X) = \frac{(b-a)^2}{12}$
	\item \textbf{泊松分布 $P(\lambda)$}：$D(X) = \lambda$
\end{enumerate}

\subsection{切比雪夫不等式}
\textbf{切比雪夫不等式（Chebyshev's Inequality）}是概率论中一个非常强大的工具，它描述了随机变量偏离其数学期望的可能性。它的魅力在于：\textbf{无论随机变量服从什么分布}（只要方差存在），它都成立。

设随机变量 $X$ 具有数学期望 $E(X) = \mu$ 和方差 $D(X) = \sigma^2$。对于任意正数 $\epsilon$，有
\begin{equation}
	P\big(|X - E(X)| \ge \epsilon\big) \le \frac{D(X)}{\epsilon^2}
	\label{eq:第四章:切比雪夫}
\end{equation}
或者其等价形式（设 $\epsilon = k\sigma$，即 $k$ 倍标准差）：
\begin{equation}
	P\big(|X - \mu| \ge k\sigma\big) \le \frac{1}{k^2}
	\label{eq:第四章:切比雪夫等价}
\end{equation}
通俗理解：随机变量 $X$ 落在期望 $\mu$ 的 $k$ 倍标准差范围之外的概率，不会超过 $1/k^2$。

\section{协方差与相关系数}
\textbf{协方差（定义）}：
\begin{equation}
	Cov(X, Y) = E\big[(X - E(X))(Y - E(Y))\big]
	\label{eq:第四章:协方差}
\end{equation}
\textbf{计算公式}：
\begin{equation}
	Cov(X, Y) = E(XY) - E(X)E(Y)
	\label{eq:第四章:协方差计算}
\end{equation}
\textbf{不相关}：若 $Cov(X, Y) = 0$，则称 $X$ 与 $Y$ 不相关。
\begin{itemize}
	\item \textbf{性质}：若 $X, Y$ 相互独立，则 $Cov(X, Y) = 0$。
\end{itemize}
\textbf{相关系数（定义）}：
\begin{equation}
	\rho_{XY} = \frac{Cov(X, Y)}{\sqrt{D(X)\,D(Y)}}
	\label{eq:第四章:相关系数}
\end{equation}
\textbf{性质}：$-1 \le \rho_{XY} \le 1$。

\textbf{$n$ 维正态分布的特性}：服从 $n$ 维正态分布的随机变量的两两不相关性与相互独立性是\textbf{等价的}。
